\subsection{モデルの分析}
モデルは作るだけでは不十分である。モデルが目的に対して十分か、矛盾していないか、正しいか、他の成果物とつながっているか、相互作用が期待どおりかを分析する必要がある。

\subsubsection{完全性の分析}
完全性とは、指定された要求がモデルに含まれ、実装や検証へつながっている度合いである。完全性は、要求獲得からコード実装まで重要である。

完全性の確認には、モデリングツールによる構造分析や状態空間到達可能性分析が使われることがある。状態空間到達可能性分析とは、モデル上の状態や経路に、正しい入力の集合から到達できるかを調べる考え方である。手作業では、レビューやインスペクションで抜け漏れを探す。

\subsubsection{整合性の分析}
整合性とは、要求、表明、制約、機能、構成要素の記述が衝突していない度合いである。たとえば、あるモデルで「利用者は退会後にログインできない」と書き、別のモデルで「退会後も注文履歴をログインして閲覧できる」と書いていれば、整合性の問題がある。

整合性確認は、ツールで自動的に行える場合がある。ただし、業務的な矛盾や意味上の矛盾は、人によるレビューが必要になることが多い。

\subsubsection{正しさの分析}
正しさとは、モデルがソフトウェア要求や設計仕様を満たし、欠陥がない度合いである。正しさには、構文上の正しさと意味上の正しさがある。

構文上の正しさは、モデリング言語の文法や記号の使い方が正しいかである。意味上の正しさは、モデルが対象領域の意味を正しく表しているかである。構文はツールで検出しやすいが、意味は関係者の知識やレビューが重要になる。

\subsubsection{追跡可能性の分析}
ソフトウェア開発では、計画書、プロセス仕様、要求、図、設計、擬似コード、コード、テストケース、テスト報告、ファイル、データなど、多数の作業成果物が作られる。これらの間には、利用する、実装する、検証する、依存する、といった関係がある。

追跡可能性とは、これらの関係をたどれることである。要求から設計、設計からコード、コードからテストへたどれると、要求が満たされていることを説明しやすい。また、変更時には影響範囲を把握しやすい。モデリングツールは、要求、設計、コード、テスト実体の間の追跡リンクを自動または手動で管理することがある。

\subsubsection{相互作用の分析}
相互作用分析は、あるタスクや機能を実現するために、モデル内の実体がどのように通信し、制御を受け渡すかを調べる。対象は、アプリケーション内部の部品だけではない。オペレーティングシステム、ミドルウェア、他アプリケーション、利用者インタフェースとの関係も含まれる。

モデリング環境によっては、シミュレーション機能を使って動的な振る舞いを確認できる。シミュレーションを一歩ずつ進めると、部品同士が意図どおり協調しているかを確認しやすい。

\par\smallskip
\begin{center}
\centering
\IfFileExists{assets/generated_figures/ch11/fig-ch11-06.png}{\includegraphics[width=0.96\linewidth]{assets/generated_figures/ch11/fig-ch11-06.png}}{\fbox{\parbox[c][48mm][c]{0.9\linewidth}{\centering 画像生成待ち\\モデル分析の五観点を示す図}}}
\par\smallskip
\refstepcounter{figure}\label{fig:ch11-06}図 \thefigure: モデル分析の五観点を示す図
\end{center}
図 \thefigure は、モデル分析の五観点を示す図を視覚的に整理したものである。
\par\smallskip
